\documentclass[12pt,UTF8]{ctexbook}

% 设置纸张信息
\input{../../../../common/page}

% 设置字体
\xeCJKsetup{AutoFallBack=true}
\setCJKfamilyfont{hei}{SimHei}
\setCJKfamilyfont{kai}{KaiTi}

% 目录 chapter 级别加点（.）
\usepackage{titletoc}
\titlecontents{chapter}[0pt]{\vspace{3mm}\bf\addvspace{2pt}\filright}{\contentspush{\thecontentslabel\hspace{0.8em}}}{}{\titlerule*[8pt]{.}\contentspage}

% 设置 part 和 chapter 标题格式
\ctexset{
	part/name= {第,卷},
	part/number={\chinese{part}},
	chapter/name={第,篇},
	chapter/number={\arabic{chapter}}
}

% 图片相关设置
\usepackage{graphicx}
\graphicspath{{Images/}}

% 列表项向右偏移
\usepackage{enumitem}
% 强制表格位置
\usepackage{float}
% 数学公式支持
\usepackage{amsmath}
\usepackage{amssymb}

\title{\heiti\zihao{0} LED}
\author{WangFei}
\date{\today}

\begin{document}

\maketitle
\tableofcontents

\frontmatter
\chapter{前言}

LED是收音机等电子设备中常用的显示和指示元件，本章节将详细介绍LED的基本特性、分类、应用以及选择原则等内容。

\mainmatter

% 增加空行
~\\

\chapter{LED的基本概念}

\section{LED的基本原理}

\subsection{LED的定义}

\subsection{LED的结构}

\subsection{LED的主要参数}

\section{LED的工作原理}

\subsection{发光原理}

\subsection{PN结的作用}

\subsection{能量转换过程}

\section{LED的正负极判断}

\subsection{直插式LED的正负极判断}

\subsubsection{引脚长度判断法}

\subsubsection{引脚形状判断法}

\subsubsection{内部结构观察法}

\subsubsection{外观标记判断法}

\subsection{贴片式LED的正负极判断}

\subsubsection{外观标记判断法}

\subsubsection{焊盘形状判断法}

\subsubsection{丝印标记判断法}

\subsection{万用表测试法}

\subsubsection{电阻测试法}

\subsubsection{电压测试法}

\subsubsection{导通测试法}

\subsection{通电测试法}

\subsubsection{低电压测试法}

\subsubsection{限流电阻测试法}

\chapter{LED的分类}

\section{按发光颜色分类}

\subsection{红色LED}

\subsection{绿色LED}

\subsection{蓝色LED}

\subsection{黄色LED}

\subsection{白色LED}

\subsection{其他颜色LED}

\section{按封装形式分类}

\subsection{直插式LED}

\subsection{贴片式LED}

\subsection{大功率LED}

\subsection{COB LED}

\subsection{其他封装形式}

\section{按功率分类}

\subsection{小功率LED}

\subsection{中功率LED}

\subsection{大功率LED}

\section{按应用分类}

\subsection{指示LED}

\subsection{显示LED}

\subsection{照明LED}

\subsection{背光LED}

\chapter{LED的特性}

\section{LED的电气特性}

\subsection{正向电压}

\subsection{正向电流}

\subsection{反向电压}

\subsection{反向电流}

\subsection{发光强度}

\subsection{发光效率}

\section{LED的光学特性}

\subsection{发光波长}

\subsection{光谱特性}

\subsection{发光角度}

\subsection{亮度}

\subsection{色温}

\subsection{显色指数}

\section{LED的热特性}

\subsection{工作温度}

\subsection{热阻}

\subsection{散热要求}

\subsection{寿命与温度的关系}

\chapter{LED在收音机中的应用}

\section{电源指示}

\subsection{电源指示灯}

\subsection{待机指示灯}

\section{调谐指示}

\subsection{调谐指示灯}

\subsection{信号强度指示}

\section{音量指示}

\subsection{音量电平指示}

\subsection{峰值指示}

\section{功能指示}

\subsection{波段指示}

\subsection{模式指示}

\subsection{故障指示}

\chapter{LED的选择原则}

\section{颜色选择}

\subsection{根据应用场景选择}

\subsection{颜色搭配考虑}

\section{亮度选择}

\subsection{根据观察距离选择}

\subsection{根据环境亮度选择}

\section{功率选择}

\subsection{根据亮度需求选择}

\subsection{根据散热条件选择}

\section{封装选择}

\subsection{根据安装空间选择}

\subsection{根据焊接方式选择}

\section{其他考虑因素}

\subsection{工作电压}

\subsection{工作电流}

\subsection{寿命要求}

\subsection{成本因素}

\chapter{LED的规格表示}

\section{型号表示}

\subsection{常见型号命名规则}

\subsection{主要厂商型号示例}

\section{参数标注}

\subsection{正向电压标注}

\subsection{正向电流标注}

\subsection{发光强度标注}

\subsection{发光波长标注}

\section{其他标注}

\subsection{极性标注}

\subsection{功率标注}

\subsection{封装标注}

\backmatter

\end{document}
